\documentclass[aps,prx,reprint,superscriptaddress,nofootinbib,longbibliography]{revtex4-2}

\usepackage{amsmath,amssymb,bm,mathtools}
\usepackage{graphicx}
\usepackage{booktabs}
\usepackage{microtype}
\usepackage{hyperref}
\usepackage{xcolor}
\usepackage{siunitx}
\usepackage{amsthm}
\usepackage{orcidlink}

\hypersetup{colorlinks=true,linkcolor=black,citecolor=black,urlcolor=blue}
\graphicspath{{figures/}}
\newtheorem{theorem}{Theorem}
\newtheorem{corollary}{Corollary}
\newcommand{\Tr}{\operatorname{Tr}}
\newcommand{\deff}{d_{\rm eff}}

\begin{document}

\title{Measurement-Accessible Quantum Tangent Geometry: Rank Baselines and Spectral Orientation}

\author{Marwan Ait Haddou\,\orcidlink{0009-0008-1734-1721}}
\email{aithaddou.marwan@outlook.com}
\affiliation{Independent Researcher}
\date{\today}

\begin{abstract}
A fixed quantum measurement can expose substantially more tangent information than a restricted observable readout retains. We study this second restriction directly. For a normalized covariance $C\succeq0$, $\Tr C=1$, of measurement-induced tangent scores in an $N$-dimensional centered score space, and a rank-$r$ readout projector $P$, we quantify retained tangent mass by $R=\Tr(PC)$. The ratio $\rho=R/(r/N)$ separates the actual retained mass from a rank-only random-orientation reference. Standard Grassmann averaging gives $\mathbb{E}\rho=1$ and
\begin{equation*}
\mathrm{Var}(\rho)=\frac{2(N-r)[N\Tr(C^2)-1]}{r(N-1)(N+2)}
\le \frac{2}{r\deff},
\qquad \deff=\frac{1}{\Tr(C^2)}.
\end{equation*}
We use this identity as a null model rather than as a new random-projection theorem. Numerically, family-balanced one- and two-body readouts remain close to the rank reference through $n=16$ even as the tangent covariance becomes strongly anisotropic. Rank matching alone is nevertheless insufficient: in generic circuits, cross-fitted leading tangent subspaces of the same rank retain substantially more mass than the physical low-weight readout. This difference is operational. At fixed circuit, measurement record, readout rank, and shot budget, changing only the readout orientation changes directional gradient energy and finite-shot signal-to-noise ratio, while random rank-matched subspaces track the rank baseline. A half-filled $U(1)$-conserving family provides a structured counterexample to generic orientation: physical low-weight $Z$ readouts are already strongly aligned with leading tangent directions over the tested finite-size range. We treat this symmetry result as a case study, not as a claim that $U(1)$ symmetry generically prevents barren plateaus or that hydrodynamics is the established mechanism. The results isolate readout orientation as a degree of freedom that is invisible to rank alone but directly controls how much measured tangent information remains usable after readout restriction.
\end{abstract}

\maketitle

\section{Introduction}

Variational quantum circuits are ultimately accessed through measurement outcomes. State-space quantities such as the quantum Fisher information (QFI), quantum geometric tensor, and natural-gradient metric characterize infinitesimal distinguishability before a particular measurement and readout are fixed \cite{BraunsteinCaves1994,Stokes2020,Haug2021}. Once a measurement is chosen, the QFI contracts to a classical Fisher geometry; once only a restricted family of observables or features is retained, the experimentally available geometry contracts again \cite{BraunsteinCaves1994,HottaOzawa2004,LuLuoOh2012}. The second contraction is the focus of this work.

This question is adjacent to, but distinct from, several established trainability results. Observable locality can change gradient concentration in parametrized quantum circuits, circuit controllability and dynamical Lie-algebra structure can alter barren-plateau behavior, and symmetry-restricted ansatzes can explore much smaller invariant sectors than generic expressive circuits \cite{Cerezo2021,UvarovBiamonte2021,Larocca2022,HolmesSharma2022}. Random measurements can also approximate quantum Fisher geometry after averaging over measurement bases \cite{LuSha2025}. These results establish that measurement choice, observable structure, and circuit architecture matter. Our question is narrower: \emph{after the quantum measurement has been fixed, and after the dimension of the retained readout has been fixed, how much tangent information is accessible to that particular readout subspace?}

The distinction matters because readout dimension does not determine readout usefulness. Consider the centered score space generated by a fixed measurement. Let $C\succeq0$, $\Tr C=1$, denote the covariance of normalized tangent-score vectors, and let $P$ be the orthogonal projector onto a retained observable subspace of rank $r$. We define
\begin{equation}
R(P,C)=\Tr(PC),
\qquad
\rho(P,C)=\frac{\Tr(PC)}{r/N}.
\label{eq:def-rho}
\end{equation}
Here $R$ is the fraction of normalized tangent mass visible to the readout. The denominator $r/N$ is the mean overlap of a rank-$r$ subspace with a trace-one covariance under random relative orientation. It is therefore a rank-only reference, not a model of a physical circuit family.

The random-subspace identity itself is standard Grassmann geometry \cite{CollinsMatsumoto2017,Bendokat2024,Fernandez2025}. We use it for a different purpose: to separate three ingredients that are otherwise easily conflated. The rank $r$ fixes the null accessibility scale $r/N$; the spectrum of $C$ controls the width of random-orientation fluctuations; and the relative orientation of $P$ and the eigenspaces of $C$ determines the actual retained mass. In particular, concentration of $\rho$ around one does not imply that $C$ is isotropic. If
\begin{equation}
\deff(C)=\frac{1}{\Tr(C^2)},
\label{eq:deff-intro}
\end{equation}
then the variance of $\rho$ is bounded by $2/(r\deff)$. Thus rank-typical overlap only requires the product $r\deff$ to be large; the ratio $\deff/N$ may still be small.

% The remainder of the manuscript follows below unchanged in the repository source.
